\scanpage{172}
\mspage{88}
% archive sheet number: 171
% lines 1-2, conclusion from previous page
la $\varinjlim_{i\in I}\Psi(i)$ existe (dans $\operatorname{Ind}'_{\pi}C$ également), et les deux sont égales.

% lines 3-11
\uline{Corollaire} \inserted{4.4.8.} Supposons que $C$ soit stable par $\varinjlim$ de cardinal $\leq\pi$ (pas seulement par celles de type $J$), \deleted{Alors \unsure{on a que}} $\underline{\operatorname{Hom}}(J,C)$ et \inserted{$I=$} $\underline{\operatorname{Hom}}(J,C)_{/\varphi}$ ont la $=$ propriété, $I$ est grand devant $\pi$, et les deux $\varinjlim_{j\in J}\varphi(j)=\varinjlim_{\unsure{j|i}\in I}\Psi(i)$ existe dans $\operatorname{Ind}'_{\pi}(C)$ (sous réserve de cofinalité des foncteurs (4.4.6.3)).

% line 12
Il suffit en effet de vérifier ceci:

% lines 13-23
\uline{Lemme} 4.4.8.1. Soit $I$ une catégorie grande devant $\pi$, et $\Psi:I\longrightarrow C$ un foncteur. Alors $\varinjlim_i\Psi(i)$ existe dans $\operatorname{Ind}'_{\pi}(C)$, et \deleted{\unsure{c'est un|est égale dans le}} \deleted{\unsure{ind-objet}} la limite inductive dans $\widehat C$ \deleted{$F$} ou encore ``c'est'' le ind-objet \inserted{grand devant $\pi$} représenté par $\Psi$. [En d'autres termes, $\operatorname{Ind}'_{\pi}(C)\subset\widehat C$ est une sous-catégorie strictement pleine \uline{stable} par $\varinjlim$ grandes devant $\pi$.]
% margin note written slanted
\margin{C'est donc cette stabilité \unsure{qui est un énoncé|est un énoncé} \unsure{nettement plus} fort! \unsure{Je} \unsure{parle prop.|par la prop.}}
